\section{2019级工科数学分析下 A卷}

\subsection{资料来源}
\begin{itemize}
  \item 试题：2019 级工科数学分析（二）A 卷无答案版，已导出页面图校正公式。
  \item 答案一：A 卷答案 1。
  \item 答案二：A 卷答案 2；两份答案有少量抽取差异，以下以题面和清晰公式为准。
\end{itemize}

\subsection{试题}

\paragraph*{试卷信息}
诚信应考，考试作弊将带来严重后果！

华南理工大学本科生期末考试 2019--2020--2 学期《工科数学分析》A卷。

注意事项：开考前请将密封线内各项信息填写清楚；所有答案请直接答在试卷上；考试形式：闭卷；本试卷共（五）大题，满分100分，考试时间120分钟。评阅教师请在试卷袋上评阅栏签名。

\paragraph*{一、填空题（共5小题，每题3分，共15分）}
\begin{enumerate}
  \item 函数 \(u=(2x^2+y^2)z^2\) 在点 \((0,1,1)\) 沿该点梯度方向的方向导数为 \underline{\hspace{5cm}}。
  \item 曲线 \(x=t+1,\ y=t^2+1,\ z=t^3+1\) 在点 \((1,1,1)\) 的切线方程为 \underline{\hspace{5cm}}。
  \item 函数 \(z=x^3+y^3-3xy\) 的极小值为 \underline{\hspace{5cm}}。
  \item 级数 \(\displaystyle \sum_{n=1}^{\infty}\frac{x^n}{1+x^{2n}}\) 的收敛域为 \underline{\hspace{5cm}}。
  \item 设 \(f(x)=\mathrm{e}^{x^2}\)，\(0\le x<\pi\)，\(S(x)\) 为 \(f(x)\) 的展成以 \(2\pi\) 为周期的正弦级数的和函数，则 \(S(0)=\) \underline{\hspace{5cm}}。
\end{enumerate}

\paragraph*{二、计算题（共5小题，每题9分，共45分）}
\begin{enumerate}
  \item 求微分方程
  \[
    y''+y'-2y=18x\mathrm{e}^x
  \]
  的通解。

  \item 设 \(z=f(x^2+y^2,x^2y^2)\)，其中 \(f\) 为任意阶可微函数，求
  \[
    \frac{\partial^2z}{\partial x^2},\qquad
    \frac{\partial^2z}{\partial y\partial x}.
  \]

  \item 计算
  \[
    \iiint_\Omega (x+y+z)\,\mathrm{d}x\mathrm{d}y\mathrm{d}z,
  \]
  其中 \(\Omega\) 是曲面 \(z=\sqrt{x^2+y^2}\) 与 \(z=\sqrt{1-x^2-y^2}\) 所围成的区域。

  \item 计算第二类曲线积分
  \[
    \oint_L xy^2\,\mathrm{d}y-x^2y\,\mathrm{d}x,
  \]
  其中 \(L:\dfrac{x^2}{a^2}+\dfrac{y^2}{b^2}=1\)，\(a,b>0\)，取逆时针方向。

  \item 计算
  \[
    I=\iint_\Sigma 2x(y-z)\,\mathrm{d}y\mathrm{d}z+(1-y^2)\,\mathrm{d}z\mathrm{d}x+(1+z^2)\,\mathrm{d}x\mathrm{d}y,
  \]
  其中 \(\Sigma\) 为柱面 \(x^2+y^2=1\) 夹在平面 \(z=0\) 和 \(z=3\) 部分的曲面，方向指向外侧。
\end{enumerate}

\paragraph*{三、解答下列各题（共2小题，每题8分，共16分）}
\begin{enumerate}
  \item 试确定 \(\alpha\) 的值，使得函数
  \[
    f(x,y)=
    \begin{cases}
      (x^2+y^2)^\alpha\sin\dfrac1{x^2+y^2}, & (x,y)\ne(0,0),\\
      0, & (x,y)=(0,0)
    \end{cases}
  \]
  在点 \((0,0)\) 处可微。
  \item 求正项级数 \(\displaystyle \sum_{n=0}^{\infty}\frac{2n+1}{n!}\) 的和。
\end{enumerate}

\paragraph*{四、证明题（共2小题，每题8分，共16分）}
\begin{enumerate}
  \item 设 \(x=x(y,z)\)，\(y=y(x,z)\)，\(z=z(x,y)\)，都是由方程 \(F(x,y,z)=0\) 所确定的具有连续偏导数的函数，证明
  \[
    \frac{\partial x}{\partial y}\cdot
    \frac{\partial y}{\partial z}\cdot
    \frac{\partial z}{\partial x}=-1.
  \]
  \item 证明函数项级数
  \[
    \sum_{n=2}^{\infty}\frac{(-1)^n(1-\mathrm{e}^{-nx})}{n^2+\sin x}
  \]
  在 \([0,+\infty)\) 上一致收敛。
\end{enumerate}

\paragraph*{五、应用题（共1题，每题8分，共8分）}
将正数 \(9\) 分成三个正数 \(x,y,z\) 之和，使得 \(u=x^2y^3z^4\) 最大。

\subsection{答案与解析}

\paragraph*{一、填空题}
\begin{enumerate}
  \item
  \[
    \nabla u=(4xz^2,2yz^2,2(2x^2+y^2)z),
  \]
  故 \(\nabla u(0,1,1)=(0,2,2)\)，沿梯度方向的方向导数为
  \[
    2\sqrt2.
  \]
  \item 该点对应 \(t=0\)，切向量为
  \[
    (x'(0),y'(0),z'(0))=(1,0,0).
  \]
  因此切线可写为
  \[
    y=1,\qquad z=1.
  \]
  \item 驻点为 \((0,0)\) 与 \((1,1)\)。在 \((1,1)\) 处 Hessian 矩阵正定，故极小值为
  \[
    1+1-3=-1.
  \]
  \item 当 \(|x|<1\) 时与 \(\sum x^n\) 同敛散；当 \(|x|>1\) 时与 \(\sum x^{-n}\) 同敛散；当 \(x=\pm1\) 时通项不趋于 \(0\)。故收敛域为
  \[
    (-\infty,-1)\cup(-1,1)\cup(1,+\infty).
  \]
  \item 正弦级数对应奇延拓，在端点 \(0\) 处取左右极限平均值：
  \[
    S(0)=0.
  \]
\end{enumerate}

\paragraph*{二、计算题}
\begin{enumerate}
  \item 特征方程为
  \[
    r^2+r-2=0,
  \]
  特征根为 \(r_1=-2,\ r_2=1\)，齐次解为 \(C_1\mathrm{e}^{-2x}+C_2\mathrm{e}^x\)。因 \(1\) 是特征根，设特解
  \[
    y^*=x\mathrm{e}^x(Ax+B).
  \]
  代入得
  \[
    6Ax+(3B+2A)=18x,
  \]
  所以 \(A=3,\ B=-2\)。通解为
  \[
    y=C_1\mathrm{e}^{-2x}+C_2\mathrm{e}^x+x(3x-2)\mathrm{e}^x.
  \]

  \item 记
  \[
    u=x^2+y^2,\qquad v=x^2y^2.
  \]
  先有
  \[
    z_x=2xf_1+2xy^2f_2.
  \]
  继续求导得
  \[
    z_{xx}=2f_1+2y^2f_2+4x^2f_{11}+8x^2y^2f_{12}+4x^2y^4f_{22},
  \]
  \[
    z_{yx}=4xyf_{11}+4xy(x^2+y^2)f_{12}+4xyf_2+4x^3y^3f_{22}.
  \]
  其中各偏导数均在 \((x^2+y^2,x^2y^2)\) 处取值。

  \item 用球坐标
  \[
    x=\rho\sin\varphi\cos\theta,\quad
    y=\rho\sin\varphi\sin\theta,\quad
    z=\rho\cos\varphi.
  \]
  两个边界分别为 \(\varphi=\pi/4\) 与 \(\rho=1\)，故
  \[
    0\le\theta\le2\pi,\qquad 0\le\varphi\le\frac{\pi}{4},\qquad 0\le\rho\le1.
  \]
  由于 \(x,y\) 关于区域对称积分为 \(0\)，只需积 \(z\)：
  \[
    I=2\pi\int_0^{\pi/4}\int_0^1 \rho\cos\varphi\cdot\rho^2\sin\varphi\,\mathrm{d}\rho\,\mathrm{d}\varphi
    =\frac{\pi}{8}.
  \]

  \item 令
  \[
    x=a\cos t,\qquad y=b\sin t,\qquad 0\le t\le2\pi.
  \]
  代入得
  \[
    \oint_L xy^2\,\mathrm{d}y-x^2y\,\mathrm{d}x
    =\frac{\pi}{4}\left(ab^3+a^3b\right).
  \]
  也可由 Green 公式化为 \(\iint_D(x^2+y^2)\,\mathrm{d}A\) 得到同一结果。

  \item 添加底面 \(\Sigma_1:z=0\) 和顶面 \(\Sigma_2:z=3\)，构成闭曲面。令
  \[
    P=2x(y-z),\quad Q=1-y^2,\quad R=1+z^2.
  \]
  因为
  \[
    P_x+Q_y+R_z=2(y-z)-2y+2z=0,
  \]
  闭曲面总通量为 \(0\)。底面外法向向下，通量为 \(-\pi\)；顶面外法向向上，通量为 \(10\pi\)。所以侧面通量为
  \[
    I=0-(-\pi)-10\pi=-9\pi.
  \]
\end{enumerate}

\paragraph*{三、解答题}
\begin{enumerate}
  \item 由
  \[
    |f(x,y)|\le (x^2+y^2)^\alpha=r^{2\alpha}
  \]
  可知若 \(\alpha>\dfrac12\)，则
  \[
    \frac{|f(x,y)-f(0,0)|}{\sqrt{x^2+y^2}}\le r^{2\alpha-1}\to0.
  \]
  此时 \(f_x(0,0)=f_y(0,0)=0\)，且函数在原点可微。若 \(\alpha\le\dfrac12\)，上述商不趋于 \(0\)，不可微。因此
  \[
    \alpha>\frac12.
  \]

  \item
  \[
    \sum_{n=0}^{\infty}\frac{2n+1}{n!}
    =2\sum_{n=0}^{\infty}\frac{n}{n!}+\sum_{n=0}^{\infty}\frac1{n!}
    =2\mathrm{e}+\mathrm{e}=3\mathrm{e}.
  \]
\end{enumerate}

\paragraph*{四、证明题}
\begin{enumerate}
  \item 由隐函数求导公式，
  \[
    \frac{\partial x}{\partial y}=-\frac{F_y}{F_x},\qquad
    \frac{\partial y}{\partial z}=-\frac{F_z}{F_y},\qquad
    \frac{\partial z}{\partial x}=-\frac{F_x}{F_z}.
  \]
  三式相乘即得
  \[
    \frac{\partial x}{\partial y}\cdot
    \frac{\partial y}{\partial z}\cdot
    \frac{\partial z}{\partial x}=-1.
  \]

  \item 对任意 \(x\in[0,+\infty)\) 且 \(n\ge2\)，有
  \[
    \left|\frac{(-1)^n(1-\mathrm{e}^{-nx})}{n^2+\sin x}\right|
    \le\frac1{n^2-1}.
  \]
  而数项级数 \(\sum_{n=2}^{\infty}\dfrac1{n^2-1}\) 收敛，由 Weierstrass 判别法，原函数项级数在 \([0,+\infty)\) 上一致收敛。
\end{enumerate}

\paragraph*{五、应用题}
在约束 \(x+y+z=9\) 下最大化 \(u=x^2y^3z^4\)。等价于最大化
\[
  \ln u=2\ln x+3\ln y+4\ln z.
  \]
由 Lagrange 乘子法，
\[
  \frac2x=\lambda,\qquad \frac3y=\lambda,\qquad \frac4z=\lambda.
  \]
故
\[
  x:y:z=2:3:4.
  \]
又 \(x+y+z=9\)，所以
\[
  x=2,\qquad y=3,\qquad z=4.
  \]
最大值为
\[
  2^2\cdot3^3\cdot4^4=27648.
\]
